\scanpage{72}
% archive foliation: 71
$\underline W^a(\mathfrak X_i\cap\mathfrak X_j)$ est \inserted{donc} \unsure{un ouvert} \inserted{à la fois aux} \deleted{\unsure{à l'ouvert}} $X_{ij}$ dans $\underline W^a(\mathfrak X_i)$ et \inserted{$X_{ji}$ dans} $\underline W^a(\mathfrak X_j)$ : cela permet de recoller les $W^a(\mathfrak X_i)=X_i$. Il faut vérifier que l'on a \deleted{$\unsure{\varphi_{ki}=\varphi_{kj}\varphi_{ji}}$} commutativité de
\diagram{\begin{tikzcd}[column sep=large]
X_{ij\unsure{k|h}}\arrow[r,"\varphi_{ji}|" description]\arrow[rr,bend left=30,"\varphi_{ki}|"]\arrow[d,equal]&
X_{ji\unsure{k|h}}\arrow[r,"\varphi_{kj}|" description]\arrow[d,equal]&X_{kij}\arrow[d,equal]\\
X_{ij}\cap X_{i\unsure{k|h}}&X_{ji}\cap X_{j\unsure{k|h}}&X_{ki}\cap X_{kj}
\end{tikzcd}}
Cela tient au fait que l'on a
\[
X_{ij\unsure{k|h}}\simeq\underline W^a(\mathfrak X_i\cap\mathfrak X_j\cap\mathfrak X_k)
\]
et que \unsure{la flèche} $\varphi_{ji}$ (p. ex.) s'identifie à l'identité.

On a ainsi défini un $X=\bigcup X_i$, $X_i\simeq\underline W^a(\mathfrak X_i)$. On définit
\[
\underline W(X)\longrightarrow\mathfrak X
\]
en notant que $W(X)=\bigcup_i W(X_i)$ \deleted{\unsure{et}} et qu'on a des \deleted{\unsure{hom}} morphismes
\[
\underline W(X_i)\longrightarrow\mathfrak X_i\subset\mathfrak X
\]
par définition de $X_i=\underline W^a(\mathfrak X_i)$. La
